\subsection{The integral symplectic model for the remaining RH characters}
\label{sec:rh-symplectic-integral-model}

The ordinary-RH specialization has a useful simplification not available
for arbitrary host characters: the set of distinguished local places is
the singleton \(\{\ell\}\). The six off-block normalizations can therefore
be made in the local ring before presenting the global ring. This avoids
a characteristic-zero elimination of integral local equations.

Retain the original datum of
Theorem~\ref{thm:rh-symplectic-arithmetic-data}, including the selected
places, \(g_0,s_0,\mathrm{ph}\), and its coefficient DVR \(O\).
Let \(k\) and \(E\) be its residue and fraction fields, respectively.
Put
\[
 \alpha=\varepsilon\omega_T^{-1}\xi,\qquad
 \Psi=\varepsilon^e\psi,\qquad \beta=\Psi/\alpha,\qquad
 \tau=\alpha\Psi .
\]
We use indices \(0,1,2,3\), the diagonal base representation
\[
 r_0=\operatorname{diag}(\tau,\Psi,\alpha,1)
\]
and the original alternating form
\[
 J=\begin{pmatrix}
 0&0&0&1\\0&0&1&0\\0&-1&0&0\\-1&0&0&0
 \end{pmatrix}.
\]
In particular the similitude is fixed to be \(\tau\).
The torus acts by conjugation by
\(\operatorname{diag}(t_1t_2,t_2,t_1,1)\). Write its four positive
root weights as
\[
 A=(1,0),\quad B=(-1,1),\quad P=(0,1),\quad T=(1,1).
\]
Their characters on the adjoint of \(r_0\) are
\(\alpha,\beta,\Psi,\tau\), respectively. The negative roots have
the inverse characters, and the weight-zero subspace of
\(\mathfrak{sp}_4\) is two-dimensional and trivial.
These are the original weights \texttt{wAl}, \texttt{wBe},
\texttt{wPs}, \texttt{wTh}, and \texttt{wt}; the three pairs
\(\pm A,\pm B,\pm T\) are the off-block roots.

\begin{theorem}[The actual local ring and its normalized component]
\label{thm:rh-symplectic-local}
There is a complete Noetherian local \(O\)-algebra \(L\), with residue
field \(k\), representing the original framed condition
\texttt{condL} (including its unnormalized spectral-projector
equations). Let \(x:L\to O\) be the diagonal augmentation, and
\(\mathfrak p=\ker x\).
There is a unique minimal prime \(Q\) of \(L\) contained in
\(\mathfrak p\). The kernel of the characteristic-zero completed
local chart at \(\mathfrak p\) is \(Q\), and some \(y\in L\) with
\(x(y)\ne0\) annihilates \(Q\). The chart is a power series ring
over \(E\) of dimension \(14\), with the following weight dimensions:
\[
 \begin{array}{c|ccc}
 \text{weight}&0&A,B,P,T&-A,-B,-P,-T\\ \hline
 \text{dimension per weight}&2&2&1 .
 \end{array}                                                   \tag{SM1}
\]
The chart can be chosen using actual entry displacements, and
has the density and saturated-power injectivity properties of
the original \texttt{IsCompl} chart. In particular \(L\) supplies
\texttt{LocBase} for this singleton local set.

Put \(L^\circ=L/Q\). The original condition \texttt{OnComp}
is represented by \(L^\circ\), on all complete Noetherian local
coefficient rings of residue field \(k\), including rings with
nilpotents and coefficient torsion.
Let \(C\) be the quotient of \(L^\circ\) by the actual off-block
entries of \(r(\mathrm{ph})\), where the blocks are
\(\{0,2\}\) and \(\{1,3\}\). Then \(C\) is a domain, represents
the original local component with that block normalization,
and its characteristic-zero completed local chart has dimensions
\[
 \begin{array}{c|ccccc}
 \text{weight}&0&A,B,T&P&-P&-A,-B,-T\\ \hline
 \text{dimension per weight}&2&1&2&1&0 .
 \end{array}                                                   \tag{SM2}
\]
In fact there is a power-series product
\[
 L^\circ \simeq C[[u_1,\ldots,u_6]],                            \tag{SM3}
\]
in which the six parameters describe the ordered pair of
Lagrangian summands determined by \(r(\mathrm{ph})\).
\end{theorem}

\begin{proof}
\emph{The literal local equations.}
The usual framed representation ring is obtained by using matrices
of finitely many generators of the quotient through which residual
lifts factor. For clarity, the kernel of the residual representation
has finitely generated maximal pro-\(\ell\) quotient, by finiteness
of its first cohomology with \(\mathbf F_\ell\)-coefficients.
The quotient of the full local group obtained this way is also
topologically finitely generated. A residual lift factors through it,
as is seen first in finite coefficient quotients and then by
separatedness. The group relations give closed equations in finitely
many matrix coordinates. Impose the fixed similitude and all the
original equations of \texttt{Ord}, not their localization or
radical. They are polynomial equations in the matrix entries
with coefficients in \(O\), and the resulting complete
Noetherian quotient is \(L\).

Explicitly, with \(I_\ell\) the absolute local inertia subgroup
and \(\chi_i=(\tau,\Psi,\alpha,1)_i\), the equations at the
distinguished element \(s_0\) prescribe
\[
 \det(X-r(s_0))=\prod_{i=0}^3(X-\chi_i(s_0)).
\]
Set \(p_i(X)=\prod_{j\ne i}(X-\chi_j(s_0))\). The other equations
are
\[
 p_j(r(s_0))r(g)p_i(r(s_0))=0\ (i<j),\qquad
 p_i(r(s_0))(r(s)-\chi_i(s))p_i(r(s_0))=0\ (s\in I_\ell).
                                                               \tag{SM4}
\]
These are precisely the polynomial projectors in the source;
none has been divided by an element of \(O\).

We describe the completed local ring at the diagonal point,
over \(E\). Its Artinian functor is the functor of continuous
lifts of the diagonal representation into finite local
Artinian \(E\)-algebras with residue field \(E\), subject
to these same equations. Here continuity is for the
finite-dimensional \(\ell\)-adic topology. To justify this
description, let \(D\) be such an Artinian algebra.
For a continuous lift, its entry displacements from \(r_0\)
have compact image in the nilpotent maximal ideal \(\mathfrak n_D\).
They lie in a finite \(O\)-lattice in that ideal. Adjoining
the finitely many products up to its nilpotence index enlarges
the lattice to one closed under multiplication. The ring
consisting of \(O\) plus this lattice is a finite complete
local \(O\)-algebra with residue field \(k\), and the lift
takes values there. Its universal map from \(L\) extends
to the characteristic-zero completion at \(\mathfrak p\).
Conversely, a map from that completion to \(D\) factors
through a finite augmentation jet. The rings \(L/\mathfrak p^s\)
are finite \(O\)-modules, since the finitely many original
entry displacements generate \(\mathfrak p\). This proves
continuity of the specialized representation in the
finite-dimensional topology, and gives the inverse
construction. No specialization through an arbitrary
discontinuous evaluation of a power series is used.

On such a coefficient algebra the eigenvalues of \(r(s_0)\)
are the distinct prescribed elements of \(E\):
with \(u=1+\ell\) they are \(u^{e+1},u^e,u,1\).
The normalized spectral idempotents, now permissible over
\(E\), give four free lines of rank one. The similitude
relation pairs a line only with its opposite line:
the product of its eigenvalue with any other non-opposite
eigenvalue differs from \(\tau(s_0)\), by an invertible
element whenever the pairing could be nonzero.
Thus the ordered flag is self-dual for \(J\).
Equations (SM4) say exactly that this flag is invariant
and that its four graded inertia characters are
\(\chi_0,\ldots,\chi_3\). Conversely any such ordered flag
satisfies (SM4), since the prescribed inertia characters
give precisely those eigenvalues at \(s_0\), and its
subspaces are the corresponding ordered sums of eigenlines.
This equivalence is over the Artinian \(E\)-algebras, not
an assertion about the characteristic-prime fibre of \(L\).

\emph{The characteristic-zero chart.}
The fixed-multiplier upper triangular group has four positive
root directions and a two-dimensional diagonal torus.
Consider a small extension of the Artinian coefficient
algebras just used. An ordered self-dual flag near the
base flag lifts: choose a symplectic basis adapted to it
and lift that basis, using smoothness of the symplectic
group. In an adapted basis, the diagonal characters
of the representation differ from the prescribed characters
by unramified characters, with two independent ratios
because the similitude is fixed. They lift across the
small extension. Indeed their Frobenius values are arbitrary
units in \(1+\mathfrak n_D\) and, by the nilpotent logarithm
and exponential, such a value defines a continuous character
of the unramified group through its \(\mathbf Z_\ell\) quotient.

Lift the upper triangular representation as a continuous
set map with these diagonal characters. Continuous lifts
exist by choosing sections in its finite-dimensional
nilpotent root coordinates. Its multiplication defect
is a two-cocycle in the small kernel, whose \(G_{\mathbf
Q_\ell}\)-action is the sum of the four positive-root
characters of the base. All four second cohomology groups
vanish. In fact local Tate duality identifies the dual
of \(H^2(E(\gamma))\) with \(H^0(E(\varepsilon/\gamma))\).
The characters \(\alpha,\beta,\tau\) are different from
\(\varepsilon\) at the chosen Frobenius lift by (BU8)
and \(\varepsilon(s_\ell)=1\). The character \(\Psi\)
is different from \(\varepsilon\) on inertia since \(e\ge2\).
Correct the multiplication defect by a continuous
one-cochain. This produces a lift with the required
ordered flag and the original equations. It proves formal
smoothness at \(\mathfrak p\).

For the tangent dimensions, the flag contributes one
direction per negative root. Each positive root has no
invariants, because its cyclotomic exponent on a subgroup
of inertia is positive. Its second cohomology is zero
as just shown. The local Euler characteristic therefore
gives dimension one for its first cohomology, and the
coboundary space has dimension one. These two dimensions
are the framed tangent directions in each positive root.
The diagonal directions are exactly the two unramified
directions. There is no further flag stabilizer dimension:
the flag associated with (SM4) is unique over the
Artinian coefficient rings. This proves (SM1).
Formal smoothness and the dimension calculation identify
\(\widehat{L_{\mathfrak p}}\) with a power series ring
in fourteen variables over \(E\).

All representation entries are equivariant for the original
two-dimensional torus, and the actual entry differences
separate framed tangents. Select from them a homogeneous
basis of the cotangent space over \(E\). The formal inverse
function theorem makes these entries the parameters of
the chart, so it is equivariant and its augmentation is
the prescribed \(x\).

For use in the integral statements, the kernel of
\(L\to\widehat{L_{\mathfrak p}}\) is the unique minimal
prime \(Q\) contained in \(\mathfrak p\): faithful
flatness of completion makes \(L_{\mathfrak p}\) a
domain, and its zero prime contracts to this unique
minimal prime. The finitely many generators of \(Q\)
are killed by a common element \(y\notin\mathfrak p\),
by the description of localization. In particular
\(x(y)\ne0\). For every \(s\ge1\),
\[
 (L/\mathfrak p^s)\otimes_O E
      \ \simeq\ \widehat{L_{\mathfrak p}}/
                       (\mathfrak p\widehat{L_{\mathfrak p}})^s .
                                                               \tag{SM5}
\]
Indeed after inverting the nonzero elements of \(O\),
every element with nonzero augmentation is already a
unit in this quotient, by the finite geometric expansion
of its inverse. Equation (SM5) proves both chart
conditions: every finite jet comes from \(L\) after
multiplication by a nonzero element of \(O\), and a
vanishing jet implies membership in the appropriate
augmentation power after multiplying by such an element
(and hence also modulo \(Q\)). These are the original
\texttt{IsCompl.dense} and \texttt{IsCompl.inj}.
The original \(L\), not just its reduced component,
represents \texttt{condL}.

We record the graded-ring convention needed here and
below. A finite set of actual entry differences can
be chosen to generate the complete ring topologically,
weight by weight, by the relative cotangent space
over \(k\) and Nakayama's lemma. Conjugation by
the torus with entries \(t_i=1+z_i\) over a formal
extension induces its formal action on the represented
functor. The kernel of this weighted presentation is
homogeneous in the coefficientwise sense. To see this,
reduce to any finite maximal-ideal jet in the target.
The weight expansion there is a finite Laurent polynomial
in the \(t_i\). Substitution \(t_i=1+z_i\) is injective
on such Laurent polynomials over any coefficient ring:
multiply by a monomial to clear negative exponents and
translate the variables in the resulting polynomial.
Thus each weight coefficient of a relation is zero
in the jet, and separatedness gives the assertion
in the complete ring. This is a weighted power-series
presentation \texttt{WPres}, and actual homogeneous
entries admit homogeneous preimages. The augmentation
sends every parameter to zero. Its algebra of finite
sums of homogeneous elements has the source's
\texttt{SmoothPt} chart: the entry parameters give
surjectivity of linear coefficients up to scalars.
For the injectivity clause, a function with augmentation
and linear part zero is in the augmentation square
after multiplication by a nonzero scalar by (SM5).
For a finite sum of homogeneous elements this is also
true in the graded subalgebra. Indeed the augmentation
ideal in the complete ring is generated by the finitely
many homogeneous parameters. Express membership in
its square by their pairwise products and project the
coefficients onto the finitely many weights needed
for the given homogeneous summands. This gives the
same expression inside the graded subalgebra. Thus
this chart verifies \texttt{IsChart.inj} with the
saturated square as specified in the source.

\emph{The actual component and its block normalization.}
The original \texttt{OnComp} condition asserts
factorization through a complete local domain carrying
a diagonal augmentation and the same local equations.
Such a domain receives a map from \(L\) whose prime
kernel is contained in \(\mathfrak p\), hence contains
\(Q\). Conversely \(L/Q\) is such a domain. This proves
the asserted functorial identification, even for
coefficient targets that are not reduced.

This component is preserved by symplectic conjugation
congruent to one, with arbitrary complete local
coefficients. One can see this directly from the
factorization definition: conjugate the representation
over \(L^\circ[[v_1,\ldots,v_{10}]]\) by the universal
formal symplectic matrix near one. The new ring is a
domain and has its diagonal augmentation at \(v=0\).
Every such conjugation is a specialization. The
same argument for the formal diagonal torus verifies
homogeneity of the component.

The element \(\mathrm{ph}\) was obtained as the
prime-to-\(\ell\) part of \(s_\ell\). For a lift
to any of our complete local rings, its matrix has
finite order prime to \(\ell\), equal to the order
of its residual matrix. In each finite maximal-ideal
quotient the limiting powers defining \(\mathrm{ph}\)
select exactly the prime-to-\(\ell\) part, and the
kernel of reduction is an \(\ell\)-group. Hence those
orders all equal the residual order; pass to the
limit using separatedness. That order divides
\(|k^\times|\), since the residual
representation is diagonal. The corresponding roots
of unity have unique Teichm\"uller lifts in \(O\).
Consequently the spectral idempotents of this matrix
are integral, including when two of its residual
eigenvalues coincide (use one idempotent per distinct
eigenvalue).

The spectra on the two prescribed blocks are disjoint.
They are the Teichm\"uller lifts of the residual
values \((\bar\alpha\bar\Psi,\bar\alpha)\) and
\((\bar\Psi,1)\) at \(\mathrm{ph}\).
Indeed the corresponding residual values at
\(s_\ell\) and \(\mathrm{ph}\) agree, and every
cross-block equality is excluded by (BU8).
Summing the spectral idempotents for the two blocks
therefore gives two canonically ordered complementary
rank-two summands, reducing to the original coordinate
blocks. Both are Lagrangian: on either block the
product of any two eigenvalues is different modulo
the maximal ideal from the similitude value, again
by (BU8). The similitude equation thus kills every
internal pairing. The cross pairing is perfect.

In the order of these two coordinate Lagrangians,
use dual bases so that the form is
\(\left(\begin{smallmatrix}0&I\\-I&0\end{smallmatrix}\right)\).
An ordered pair of complementary Lagrangians near
the coordinate pair is uniquely a pair of graphs
of symmetric \(2\)-by-\(2\) matrices \(U,V\) with
entries in the maximal ideal. The graphs are complementary
since \(I-UV\) is invertible. The canonical matrix
\[
 h=\begin{pmatrix}I&V\\U&I\end{pmatrix}
                \begin{pmatrix}I&0\\0&(I-UV)^{-1}\end{pmatrix}
                                                               \tag{SM6}
\]
is symplectic, congruent to one, and carries the
coordinate pair to these graphs. This follows
by direct multiplication with the displayed alternating
form, using the symmetry of \(U,V\).
Conjugating by \(h^{-1}\) normalizes the two blocks
of the matrix at \(\mathrm{ph}\), without changing
the component. Conversely conjugating any such
normalized representation by \(h\) recovers a
representation with the given pair of graphs.
The constructions are inverse and are natural in
all the complete local coefficient rings. This
proves (SM3), with the normalized quotient \(C\)
defined by the actual off-block entry equations.
In particular \(C\) embeds into the domain \(L^\circ\)
and is a domain; no integral saturation is involved.

At the characteristic-zero base, the graph parameters
remove precisely one direction in each of the six
off-block root spaces. This is seen either from
the graphs or by infinitesimal conjugation:
on each such root the commutator with
\(r_0(\mathrm{ph})\) is nonzero by the disjoint
spectra. The split product (SM3) and the smooth
chart of \(L^\circ\) give a regular chart of \(C\)
over \(E\); it has precisely the dimensions (SM2).
Choosing actual entry differences for its tangent
coordinates and for its weighted complete presentation
works as above. The same finite-jet argument proves
its saturated-square injectivity. This completes the
local constructions.
\end{proof}

\begin{theorem}[The actual relative equations and integral rank-four model]
\label{thm:rh-symplectic-model}
For these RH data the original \texttt{Fr} relative relation
counts hold over \(L\). More generally they hold over the
normalized component \(C\) of Theorem~\ref{thm:rh-symplectic-local},
with the local representation fixed over that base.
In either case there are \(n_j\ge0\) new variables of weight \(j\),
chosen from actual global entry displacements, with
\[
 \begin{split}
 n_j&=0\quad(j\notin\{0,P,-P\}),\\
 b_j&=n_j+2\quad(j=0),\\
 b_j&=n_j+1\quad(j=P,-P),\\
 b_j&=0\quad(j\notin\{0,P,-P\}),
 \end{split}                                                   \tag{SM7}
\]
where \(b_j\) is the number of homogeneous equations of weight \(j\)
in a presentation relative to the given local ring; zero equations
can be added to attain these numbers.

In particular an actual \texttt{Model} exists, with the full
original integral relation ideal and finite-modulus derivation
condition. One may take its ambient domain to be
\[
 R=C[[X_{j,i}: j\in\{0,P,-P\},\ 1\le i\le n_j]],
                                                               \tag{SM8}
\]
with augmentation zero on the new variables, and
\[
 r=n_{-P}+1,\qquad r'=n_P+2,\qquad m=n_0+2.                    \tag{SM9}
\]
Its universal normalized quotient has exactly the prescribed
\texttt{IsPt} functor, rather than just a chosen set of
characteristic-zero points. Its equation lists are those
of \texttt{Model}: \(r+1\) equations of weight \(-P\),
\(r'\) of weight \(P\), and \(m\) of weight zero.
\end{theorem}

\begin{proof}
\emph{Relative cohomology.}
Let \(\Sigma\) be the set of primes dividing \(2\ell N\);
use \(G_{\mathbf Q,\Sigma}\) and the original maps from the
absolute local decomposition groups. For \(q\in\Sigma\setminus
\{\ell\}\), put \(N_q=I_q\cap\ker\bar r_0\).
Since all the original finite character values on such
inertia have order prime to \(\ell\), reduction is faithful
on them. Thus the base representation kills \(N_q\),
and \(I_q/N_q\) has order prime to \(\ell\).
In particular the local condition of factoring through
\(D_q/N_q\) is formally smooth with fixed similitude:
this group has \(\ell\)-cohomological dimension at most one,
by its finite prime-to-\(\ell\) inertia and procyclic
unramified quotient. Lifting into the symplectic group
is unobstructed over each small extension, with the
prescribed similitude, by that second-cohomology vanishing.

Let \(M_j\) be the \(j\)-weight subspace of
\(\mathfrak{sp}_4(k)\), with the action through \(\bar r_0\).
Use the same relative cochain complex as in (RP5):
\[
 {\cal C}_j=
 \operatorname{Cone}\left(
 C^\bullet(G_{\mathbf Q,\Sigma},M_j)\longrightarrow
 C^\bullet(D_\ell,M_j)\ \oplus\
 \bigoplus_{q\ne\ell}
 \bigl(C^\bullet(D_q,M_j)/C^\bullet(D_q/N_q,M_j)\bigr)
 \right)[-1].                                                 \tag{SM10}
\]
All cochains are continuous, and the full degree-zero term
at \(\ell\) is retained. It is the relative complex for
lifts with the local representation fixed as a matrix
representation (not fixed only up to conjugacy).
Its first cohomology identifies with global cocycles
vanishing at \(D_\ell\) and on each \(N_q\):
subtract the global coboundary from the local degree-zero
component at \(\ell\) to make that component zero.
This representative is unique, since the relative
degree-zero map at \(\ell\) is the identity on the
coefficient module.

For \(j=\pm A,\pm B,\pm T\), the original six \texttt{XVan}
assertions proved in Theorem~\ref{thm:rh-symplectic-arithmetic-data}
make every such cocycle a global coboundary.
But \(H^0(D_\ell,M_j)=0\) by (BU8), so the requirement that
it vanish at \(D_\ell\) kills its coboundary parameter.
Consequently \(H^1({\cal C}_j)=0\) on the six off-block
roots. The only other nonzero coefficient modules are
\(M_{P},M_{-P},M_0\).

For all these modules, \(H^0(\mathbf Q,M_j^\vee(1))=0\).
The weight-zero and the two \(\Psi\) directions have even
character, whereas the residual cyclotomic character is odd.
The six odd root characters have nontrivial dyadic ramification,
and so cannot equal the residual cyclotomic character, whose
conductor divides \(\ell\). Poitou--Tate duality therefore
gives the surjectivity of restriction in global degree two
to the sum of the local degree-two groups
\cite{AuxTowerDuality}. Since \(D_q/N_q\) has no cohomology
above degree one, this implies \(H^3({\cal C}_j)=0\).
The restriction at \(\ell\) gives \(H^0({\cal C}_j)=0\).
All cohomology spaces here are finite-dimensional, by
the finite global and local Galois cohomology theorems.
The quotient complexes at \(q\ne\ell\) have Euler
characteristic zero: the local Euler characteristic is
zero, and \(H^1(D_q/N_q,M_j)\) has the same dimension
as \(H^0(D_q/N_q,M_j)=H^0(D_q,M_j)\), by the
procyclic quotient calculation.
The local Euler characteristic at \(\ell\) is
\(-\dim M_j\), and the global one over \(\mathbf Q\)
is \(-\dim M_j^{c=-1}\). Therefore
\[
 \dim H^2({\cal C}_j)-\dim H^1({\cal C}_j)
   =\dim M_j^{c=1}
   =\begin{cases}
      2& j=0,\\
      1& j=P,-P,\\
      0& \text{otherwise}.
     \end{cases}                                             \tag{SM11}
\]
This is the original \texttt{Raw.Nums} calculation with
\(S_b=\varnothing\); no symmetry assumption on the
sizes of the two \(\Psi\)-cohomology groups is needed.

\emph{The integral relative presentation.}
Take the local ring to be either \(L\) or \(C\), and
require the actual local matrix representation from
that ring at \(D_\ell\). Require the fixed similitude
globally and the literal kernels \(N_q\) elsewhere.
The resulting complete local global ring \(U_{\mathrm{gl}}\)
is represented by finitely many entry variables, as
in the local construction. The local map is the specified
map from the original local ring; it is not assumed flat.
For the original \texttt{Fr} the only local framing
matrix is the identity, because the set of distinguished
places is \(\{\ell\}\). Thus the construction over \(L\)
is exactly that framed functor.

The tangent and obstruction argument of
Theorem~\ref{thm:support-relative-presentation} applies
to (SM10). More explicitly, relative framed tangents
are the cocycles just described, hence choose
\(n_j=\dim H^1({\cal C}_j)\) new entry coordinates
in each weight. The weighted power-series map from
the local ring with those variables onto
\(U_{\mathrm{gl}}\) is surjective by complete
Nakayama, and its kernel is coefficientwise homogeneous
by the formal torus argument. To bound its relations
one applies the small-extension obstruction map to
the dual of the minimal relative relation space.
A global matrix lift can be chosen with the prescribed
similitude, since that group is smooth; at \(\ell\)
compare with the already specified local matrix lift,
and at other primes with the unobstructed lifts
through \(D_q/N_q\). Their multiplication defects and
differences give precisely the degree-two cocycle of
(SM10). Vanishing of this obstruction gives a global
lift with the same local base map. A section of the
corresponding pushout of the relation extension
would make the original relation functional zero:
changing a new variable by an element of the small
kernel does not change the image of any relation,
since all its linear coefficients belong to the
maximal ideal by minimality of the relative variables.
This is the integral relation-space injection proved
in the earlier relative presentation theorem.
It gives at most \(\dim H^2({\cal C}_j)\) equations
in weight \(j\); the same weight is obtained because
the matrix entries of the cocycle with values in a
weight-\(j\) small kernel have that weight. Padding
with zero equations and using (SM11) gives (SM7).
If a weighted power-series presentation over \(O\)
for the chosen local base is specified, retain every
one of its equations and lift these new relative
equations homogeneously. No local equation or
coefficient torsion in the global quotient is
discarded. This proves both relative assertions,
including the original \texttt{Fr} counts.

\emph{The actual ambient domain and quotient.}
Now use \(C\) as the local base. Its local matrix
representation already satisfies the block
normalization at \(\mathrm{ph}\), including on
nilpotent coefficient rings, so the global functor
over it has precisely this normalization as well.
Its ambient ring is (SM8), a domain. The remaining
two original normalizations are
\[
 \bigl(r(g_0)+\tau(g_0)r(g_0)^{-1}\bigr)_{02}=0,\qquad
 \bigl(r(g_0)+\tau(g_0)r(g_0)^{-1}\bigr)_{20}=0 .
                                                               \tag{SM12}
\]
They have weights \(P,-P\), respectively, and vanish
at \(r_0\). Add their actual homogeneous preimages
to the relative relation ideal. This produces the
complete quotient \(S\) of \(R\) and its representation
\(\rho_u\). By the functorial characterization of
\(C\) and the explicit equations (SM12), the
points of \(S\) are exactly the source's
\texttt{IsPt}: they obey the fixed similitude,
\texttt{UnrT}, the block and the two marked conditions,
and the literal local \texttt{OnComp} condition.
There are no other local component conditions for
the present \(S_b=\varnothing\).
Maps from \(S\) are determined by the matrix entries:
the local base map is uniquely determined by restriction
of the matrix representation, and the new variables
are chosen actual entry displacements. This is
\texttt{Lift.Generates} for this quotient.

The weighted chart of \(R\) is the one of \(C\)
with the additional variables freely adjoined.
By (SM2) and (SM7) its weight multiplicities are
one each at \(A,B,T\), \(n_{-P}+1\) at \(-P\),
\(n_P+2\) at \(P\), and \(n_0+2\) at zero,
and vanish elsewhere. The finite-jet and homogeneous
augmentation-square argument in the local theorem
gives the exact \texttt{SmoothPt} assertion for the
graded subalgebra of \(R\). This is the source's
\texttt{tDims} with (SM9).
The number of equations at \(-P\), after (SM12),
is \(n_{-P}+2=r+1\), at \(P\) is \(n_P+2=r'\),
and at zero is \(n_0+2=m\). There are no equations
in the six off-block roots. Every zero-weight
equation vanishes at the diagonal augmentation.
Also \(r\ge1\) by (SM9).
Choose the corresponding lists \(F,E_+,Z\);
their ideal is the entire kernel of the quotient
map \(R\to S\). All entries of \(\rho_u\) lift
homogeneously to \(R\), and any such lift evaluates
at the augmentation to the corresponding entry
of \(r_0\), since the entire relation ideal vanishes
at that augmentation. These verify the chart,
equations, homogeneity and augmentation fields of
\texttt{Model}.

\emph{The finite-modulus derivation.}
It remains important to check the actual integral
lifting field of that structure. Let \(0\ne a\in O\)
and let \(z\) be any original \texttt{ShaCoc}\((a,z)\):
it is an \(O/(a)\)-valued continuous cocycle for
\(\Psi^{-1}\), with the original inertia-kernel
condition, trivial local class at \(D_\ell\), and
\(z(g_0)=0\). In the matrix convention used here
the root vector is exactly
\[
 X_-=E_{20}+E_{31}\in\mathfrak{sp}_4(O)_{-P}.
\]
Set \(D_a=O\oplus(O/(a))\delta\), with \(\delta^2=0\).
The matrices
\[
 r_a(g)=(1+\delta z(g)X_-)r_0(g)                              \tag{SM13}
\]
form a representation with the fixed similitude:
the root action is \(\Psi^{-1}\) and the
linear symplectic equation follows from
\(X_-^{\mathsf t}J+JX_-=0\).
They kill the prescribed inertia kernels outside
\(\ell\) by the defining condition on \(z\).

Write \(z(g)=(\Psi(g)^{-1}-1)m_a\) on \(D_\ell\).
Locally (SM13) is the conjugation of \(r_0\) by
\(1-\delta m_aX_-\). This is an integral symplectic
conjugation in the two blocks, congruent to one.
The representation therefore factors through \(C\)
at the local place: it satisfies the component
condition and its matrix at \(\mathrm{ph}\) is
block diagonal. For the two global equations (SM12),
\(z(g_0)=0\) also implies \(z(g_0^{-1})=0\).
Hence (SM13) satisfies the entire original
\texttt{IsPt} functor, and gives a map \(S\to D_a\)
reducing to the diagonal augmentation modulo \(\delta\).

The induced modular derivation on \(C\) has an actual
integral lift: choose any \(m\in O\) lifting \(m_a\),
and conjugate \(r_0\) locally by \(1-\delta mX_-\)
over \(O[\delta]/(\delta^2)\).
This stays in the original component and the block
normalization for exactly the same reasons. The
corresponding derivative \(d_C\) is supported in
weight \(-P\). Lift the modular derivative of each
new variable of weight \(-P\) to \(O\), and set the
derivatives of all other new variables to zero.
By the free power-series description (SM8), these
values and \(d_C\) give an integral \(O\)-derivation
\(d_R\) at the augmentation, of weight \(-P\).
It reduces modulo \(a\) to the derivative of the
actual map \(R\to S\to D_a\). Consequently it sends
the entire integral relation ideal into \((a)\),
not merely a list of equations or their
characteristic-zero localizations.

Finally the two entries specified by \texttt{Model.der}
are exactly
\[
 \begin{split}
 d_R(f)\bmod a&=\tau(g)z(g)
       &&\text{if }f\mapsto(\rho_u(g))_{20},\\
 d_R(f)\bmod a&=\Psi(g)z(g)
       &&\text{if }f\mapsto(\rho_u(g))_{31}.
 \end{split}                                                   \tag{SM14}
\]
They follow from (SM13), independently of the chosen
preimages \(f\), since the whole kernel maps into
\((a)\). Thus the required derivation lifting is
valid for every nonzero modulus and every original
normalized cocycle, with precisely the prescribed
coefficients and signs. This finishes the construction
of the actual rank-four integral model.
\end{proof}

The theorem supplies the model input of
\texttt{SymplecticGSp4.GaloisRepIneqOfAllIrreducibleNonvanishing}
for the remaining ordinary-RH Case II family. The subsequent
Theorem~\ref{thm:rh-sp4-conditional-length} supplies the integral
length comparison from weight-zero finiteness, block-endpoint
exclusion and fixed-field cuspidal attachment. Those three actual
arithmetic premises remain open; the rigidity implication and
uniform trace count follow under the stated attachment premise. In particular neither the characteristic-prime
fibre nor the varying Selmer lengths are controlled just
by the dimension of the characteristic-zero chart.
